% =====================================================================
%  Harald Høffding, Om Religionsfrihed og nogle dermed forbundne psykologiske og etiske Spørgsmaal
%  In: Tilskueren. Maanedsskrift for Literatur, Samfundsspørgsmaal og almenfattelige videnskabelige Skildringer,
%  udgivet af N. Neergaard.  Tredje Aargang.  Kjøbenhavn: P. G. Philipsens Forlag; Thieles Bogtrykkeri.  Juni--Juli 1886, pp. 448-459.
%  Antiqua.
%
%  TRANSCRIPTION -- GENERATED; DO NOT EDIT BY HAND.  Høffding's four pieces in Tilskueren 1886-1890 (Om Religionsfrihed;
%  the three contributions to the Nietzsche debate with Georg Brandes) are built by ONE pipeline
%  (ebook-method/tilskueren-1886-90/) into a combined working file, which split.py cuts into four transcriptions.
%  Rebuild from the repo root:   sh ebook-method/tilskueren-1886-90/rebuild.sh
%
%  STATUS: COMPLETE (2026-10-15): the volume's title page and the whole text, pp. 448-459, with its notes.
%
%  EDITION.  The only printing in Tilskueren.  Text of record: Google's scan of the Harvard copy, Internet Archive
%  tilskuerenmaaned1886nngoog, 600 ppi grey JP2 leaves with the Archive's ABBYY OCR; extract ~/bibliotek/Høffding, Harald/1886-tilskueren-3-pp448-459.pdf (SCANS.tsv, sha256),
%  made by ia_extract.py (title leaf, Indhold leaves, the article).  Page range confirmed from the volume's Indhold and the
%  signature.  Not reprinted in Høffding's Mindre Arbejder.  No e-text found (henrikpontoppidan.dk has only the Nietzsche-debate pieces).  The note on the title (a lecture in Studenterforeningen, 20 March 1886) is the print's.
%
%  ACCURACY (2026-10-15; one collation for all four texts, 76 pp., ~29,000 words).
%  (a) The DRAFT is the Archive's ABBYY OCR (the layer of the working PDF), turned into LaTeX by draft.py; it was collated
%      word by word with an independent tesseract reading (dan, tessdata_best) of the same images (collate/check.py,
%      witness 'ocr'): 650 disagreements.  check.py's own rules settled 44.
%  (b) A THIRD reading voted: Project Runeberg's OCR of the same Google images (never proofread there), by context
%      (rules.py umn-*, 372 decisions), guarded so that a vote never makes a word the corpus does not attest; other
%      rules (the layer's glyph faults against an attested word, accents tesseract adds, the ø rule, the house thin space)
%      settled 70.  The vote was AUDITED at the image: 40 rule decisions sampled blind and 30 more suspects read at zooms:
%      7 Runeberg-for-the-draft votes were wrong (ABBYY and Runeberg sharing a fault: bg/og, kim/kun, bar/har,
%      Anstrengelser/Anstrængelser, Vist/vist, én/en, h Vad/hvad), all mended; the ~30 letter-changing votes for the draft
%      not looked at are tesseract noise against a real word.
%  (c) The rest were read at the image BLIND (round b1: 303 crops on 31 sheets, six readers, 62/62 planted controls right):
%      203 open checks, 23 places where Runeberg disputed both OCRs, 54 places where the henrikpontoppidan.dk e-text
%      differs (a FINDING AID only: someone else's transcription, never a vote -- it normalises the print's spelling, e.g.
%      Nietzsche for the 1889 print's Nietsche), and the audit; ~150 decided in the main session from the readers' lines,
%      ~25 at 300-dpi zooms.  OPEN 0.  An unattested-word sweep of the collated text (353 words in no other transcription)
%      was read through; its doubtful words were looked at in zooms (one misprint found: forkellige).
%  (d) STRUCTURE read at the image (reader K): the title pages, the article heads, the section numerals (two the draft
%      had lost), the verse, the divider; every FOOTNOTE read whole (reader N: 11 notes; compared word by word with the
%      collated notes: 0 word differences, the reader right on the closing quotes).
%  (e) LETTERSPACING (the print's emphasis, set \emph{} here) is MEASURED at the image, not read: two blind line rounds
%      (l1, 256 lines; l2, the 134 doubtful ones again) agreed on only 40 of 134 lines, so spacemeasure.py decides by the
%      white between the letters of each word in tesseract's boxes (bimodal: 1-4 px roman, 6-8 px spaced; threshold 5 px;
%      isolated words must have every gap wide).  Checked at zooms on 16 words where the measure and the e-text's <em>
%      differ: the measure right on 11; three of its false positives (heavy isolated words: Ævne, frem, Side) led to the
%      isolation rule; three misses mended by hand after the page check (SPACEFIX in spaceapply.py).
%  (f) SAMPLED PAGE CHECK (pagecheck.py, seed 20261015; two blind readers): 6 whole pages, one or two from each piece
%      (Religionsfrihed p. 451; Demokratisk Radikalisme pp. 852, 868; Gensvar pp. 132, 140; Epilog p. 343), 2202 words:
%      1 WORD error (Talemaader for the print's Talemaaler: a correction wrongly taken from the e-text finding aid),
%      1 PUNCT (a speck read as '), both mended; LETTERSPACING 15 differences: at 300-dpi zooms 5 transcription errors
%      (tre, Kopier twice, Genmæle, fortsatte; mended), 4 where the transcription is right, 6 not decidable at 300 dpi
%      (the ends of runs in wide-set lines).  Word error rate after mending 0 in 2202 (95% upper bound ~1.4 per 1000);
%      letterspacing remains the least certain feature: ~2-3 per 1000 words.
%  PRINTED AS IS (misprints kept, marked at the line):
%    p. 459: „ich“ Goethe: „Gefühl ist alles“; zoom z1, reader K
%  CONVENTIONS.  Page turns: % --- p. N --- and \opage{N} (the printed page); a word divided over a page turn is split with
%  %.  Footnotes as \footnote at their marks, numbered per page as printed (1)).  \emph{} = letterspacing; signatures in
%  small capitals as printed.  The running heads and the sheet signatures are not transcribed.  The ø rule (CLAUDE.md).
%  Quotation marks as printed: »...« and „...“ both occur, and the print sometimes doubles them (»»Zunächst).
% =====================================================================
\documentclass[12pt,a4paper,openany]{book}
\usepackage[utf8]{inputenc}
\usepackage[T1]{fontenc}
\usepackage{libertinus}
\usepackage{libertinust1math}
\usepackage[danish]{babel}
\usepackage[protrusion=true,expansion=true]{microtype}
\usepackage{setspace}
\usepackage[margin=1.3in]{geometry}
\usepackage{fancyhdr}
\usepackage{hyperref}
\hypersetup{pdftitle={Om Religionsfrihed}, pdfauthor={Harald Høffding}, hidelinks}
\pagestyle{fancy}
\fancyhf{}
\fancyhead[LE,RO]{\thepage}
\fancyhead[C]{\textit{H. Høffding: Om Religionsfrihed (1886)}}
\fancypagestyle{plain}{\fancyhf{}\fancyhead[LE,RO]{\thepage}}
\setstretch{1.3}
\usepackage{marginnote}
\usepackage{graphicx}
\DeclareUnicodeCharacter{0254}{\reflectbox{c}}   % ɔ (id est)
\renewcommand{\thefootnote}{\arabic{footnote})}   % the print numbers each page's notes 1), 2)
\newcommand{\opage}[1]{\marginnote{\footnotesize\textit{[#1]}}}
\newcommand{\sectionnum}[1]{\par\medskip\begin{center}#1\end{center}\smallskip}
\newcommand{\asterism}{\par\begin{center}*\quad*\\[-0.4em]*\end{center}}
\newcommand{\signature}[1]{\par\smallskip\hfill #1\hspace*{2em}\par}
\pdfstringdefDisableCommands{\renewcommand{\footnote}[2][]{}\renewcommand{\opage}[1]{}}

\begin{document}

\begin{titlepage}
\begin{center}
\vspace*{2em}
{\LARGE\bfseries \emph{TILSKUEREN.}}\\[1em]
{\large \emph{MAANEDSSKRIFT}}\\[0.5em]
{\small FOR}\\[0.5em]
LITERATUR, SAMFUNDSSPØRGSMAAL OG ALMENFATTELIGE\\
VIDENSKABELIGE SKILDRINGER.\\[1.5em]
{\small UDGIVET AF}\\[0.5em]
{\bfseries N. NEERGAARD.}\\[1.5em]
\emph{TREDJE AARGANG.}\\
\vfill
KJØBENHAVN 1886.\\[0.4em]
{\bfseries P. G. PHILIPSENS FORLAG.}\\[0.3em]
{\small THIELES BOGTRYKKERI.}
\end{center}
\end{titlepage}
\clearpage

% --- p. 448 ---
\setcounter{footnote}{0}%
\opage{448}
\begin{center}{\large\bfseries Om Religionsfrihed og nogle dermed forbundne psykologiske\\
og etiske Spørgsmaal\footnote[1]{I det væsentlige et Foredrag, holdt i Studenterforeningen den 20de Marts 1886.}.}\\[0.5em]\rule{4em}{0.4pt}\end{center}\medskip

For det meste forstaar man Religionsfrihed som noget blot
\emph{ydre} og noget blot negativt, nemlig det, som vor Grundlov udtrykker
saaledes, at Ingen paa Grund af sin Trosbekendelse kan
berøves Adgang til den fulde Nydelse af borgerlige eller politiske
Rettigheder. Allerede dette er noget meget betydningsfuldt, naar
man erindrer, hvilken Modstand det deri udtalte Princip har haft
at kæmpe imod, før det trængte igennem. Men for den, der
gennemtænker den Forandring, som med Anerkendelsen af dette
Princip er indledet i det menneskelige Livs hele Aand og Retning,
--- for den, der mener, at det ydre og det indre vel kan træde i
en vis Modsætning til hinanden, men dog i Længden maa svare
til hinanden, hvis Livet skal være sundt og kraftigt, --- for den,
der ser længere end til det blot udvortes og foreløbige, --- for den
ligger der i Religionsfriheden psykologiske og etiske Problemer af
den største Betydning, Problemer, som ingen, der tager Livet
alvorlig, kan undgaa at blive stillet overfor. Jeg vil her forsøge
at fremdrage og belyse nogle af disse Problemer. For at de ret
kunne træde frem, maa jeg begynde med nogle historiske, eller
som man nu plejer at sige, sociologiske Bemærkninger.

\sectionnum{I.}

Positiv Religion er \emph{social} Religion, idet den bygger paa
Tradition, paa Ordets Forkyndelse. Under forskellige Former
genfindes dette Hovedtræk i alle Religioner. Det er ikke de
% --- p. 449 ---
\setcounter{footnote}{0}%
\opage{449}enkelte Individer, der udfinde deres Religion, men de indvies i
den, der lever i Slægten. Saa snart man tillægger det enkelte Individ
--- dets Trang, dets Følelse, dets Dømmekraft --- en absolut
Afgørelsesret i religiøs Henseende, er man allerede inde paa Rationalismen.
Religionen er en \emph{historisk} Magt; det er Forkyndelsen
af en hellig Historie, som danner det Grundlag, paa
hvilket der bygges.

I sin simpleste Form bygger Religionen paa en Tradition,
der indenfor én og samme Familje forplantes fra Slægt til Slægt.
De afdødes Aander dyrkes ved Ofre og Gaver, Arneilden holdes
ved lige til deres Ære, og deres Bud agtes som højeste Lov. I
denne \emph{Familjereligions} Udøvelse deltage kun de, der høre til
Slægten; Familjefaderen, som leder den, staar som den højeste
levende Myndighed i religiøs som i andre Henseender. Endnu
den Dag i Dag er denne Art Kultus den mest udbredte paa
Jorden, idet den bestaar under, ved Siden af eller som Led af de
nationale Religioner. Kun i Kristendommen og i Muhammedanismen
er denne Dyrkelse helt forsvunden\footnote[1]{\emph{Henry Maine}: Early Law and Custom. p. 57 ff. --- Det var \emph{Fustel de Coulange}, som først i sin interessante Bog „La cité antique“ paaviste den store Betydning og primitive Karakter af Familjereligionen indenfor den græske og romerske Religion.}.

\emph{Den nationale Religion} opstod ved, at hele Staten fik
sine fælles Guder og Helte, sin fælles Arne. Staten var ligesom
Familjen oprindelig baade et religiøst og et politisk Samfund.
Der var ingen Forskel paa Stat og Kirke, saa lidt som paa religiøse,
juridiske og moralske Love. Det var enhver Borgers Pligt
at deltage i den offentlige Kultus, men fremmede vare udelukkede
fra den. Hver Stat havde sine Guder, ligesom hver Familje
havde sine.

Et betydningsfuldt Vendepunkt indtraadte, da \emph{Ideen om
Kirken som universelt religiøst Samfund} opstod. Familjeforskelligheder
 og nationale Skranker skulde nu kun have underordnet
Gyldighed. Der skal ikke mere være Jøde eller Hedning,
Græker eller Barbar. Et rent aandeligt Samfund skal eksistere.
Men ogsaa her er det religiøse Samfund, den aandelige Enhed betinget
ved den fælles Overlevering af en hellig Historie, fra hvilken
de enkelte hente den højeste Lov for deres Liv og den dybeste
Betryggelse for deres Skæbne i denne og i en hinsidig Verden.
% --- p. 450 ---
\setcounter{footnote}{0}%
\opage{450}For den europæiske Menneskeheds Vedkommende gjorde Kristendommen
denne Ide gældende paa en sejerrig Maade.

I Begyndelsen stillede Kirken sig i bestemt Sondring fra
andre menneskelige Samfund som et Rige, der ikke var af denne
Verden. Den levede i Forventningen om, at denne Verden vilde
ophøre. Da denne Forventning ikke opfyldtes, begyndte Kirken
sin kulturhistoriske Mission. Den paatrykte snart det indre og det
ydre Liv sit Præg, hvad enten dens Herredømme fremtraadte derved,
at den ligefrem herskede over Staten, eller derved, at Staten
i den saa' et uundværligt Redskab til at sikre Fred og Lydighed.
Efter den middelalderlige Kirkestat, hvor Paven var Kejserens
Overhoved, fulgte det »territoriale« Princips Herredømme, hvor
Landsherrens Religion ogsaa skulde være Undersaatternes (cujus
regio, ejus religio).

Da kommer fra det 18de Aarhundrede af Principet om
Statens religiøse Neutralitet frem. Hvad enkelte isolerede Tænkere
havde fordret, proklameredes nu som Lov for store Nationer. De
nordamerikanske Fristaters Forfatning er den første, som er bygget
paa Religionsfrihedens Princip.

Et stort Gode er dermed opnaaet. Den aandelige Kamp
kan nu føres friere, da de ydre Vilkaar ere mere lige for alle.
Saa mægtige Fordomme og Fanatisme endnu stadig kunne være,
saa raade de dog ikke mere over saadanne Midler som Baal og
Brand eller udvortes Forfølgelse.

Men et stort Ideal, en stor Tanke er opgiven: Tanken om
en fælles Tro, om inderligt aandeligt Fællesskab for den hele
Menneskeslægt. Og en stor Kraft synes at maatte gaa tabt, den,
der ligger i Sammenslutning formedelst fælles Anskuelse om Livets
største Problemer. Vi have nu Lov hver at blive salige paa sin
Façon. Men vil det ikke let føre med sig, at vi sejle hver sin Sø
og tabe den gensidige Forstaaelse?

\sectionnum{II.}

Man vil fra kirkelig Side svare, at det har ingen Fare: Opløsningen
vil kun være forbigaaende, og paa Frihedens Grund vil
en ny Tros-Enhed, sandere end den tidligere tvungne Enhed, opnaas.
Og i Fritænkernes Lejr haabe mange paa, at den Tid ikke
er fjærn, da alle fornuftige, ærlige og dannede Mennesker have
vendt Kirken Ryggen, saa at i det højeste kun Dummerhoveder og
Hyklere blive tilbage.

% --- p. 451 ---
\setcounter{footnote}{0}%
\opage{451}Jeg deler ingen af disse Forventninger. --- Lige saa lidt som
den katolske Kirke trods sit store Haab og sine ivrige Anstrængelser
vil kunne tilbageerobre Protestanterne, lige saa lidt vil
Kirken kunne tilbageerobre Fritænkerne. Lige saa lidt som der i
Katolicismen er Mulighed for et saadant Opsving eller en saadan
Udfoldelse af aandelig Overlegenhed, at Protestanterne derved
kunde bevæges til at opgive deres Trosfrihed, deres Ret til selvstændig
religiøs Prøvelse, lige saa lidt er der nogen Mulighed for,
at Kirken skulde bringe de Motiver, som gør Folk til Fritænkere,
ud af Verden. Den middelalderlige Kirke formaaede endnu at
samle al Tidens aandelige Kraft i sig. Men Livet er nu for rigt,
for selvstændigt, for mangfoldigt, til at noget saadant er muligt.
De store Problemer ere blevne stillede saa klart og skarpt, at
ingen dogmatisk Formel kan dække over dem. Trangen til, ---
hvad enten man mener at der gives en Løsning af Livets Gaader
eller ikke, --- dog at behandle dem paa sin egen Maade er bleven
saa mægtig, at den ikke igen vil kunne kues. Den kan dysses
hen for en Tid, men vil stedse vaagne paany, saa længe der bestaar
et højere aandeligt Liv. Og hvad der gælder for Erkendelsens
Vedkommende, gælder ogsaa for Fantasiens, Følelsens og
Viljens Liv. Ogsaa her er den frie Rørelse af Kræfterne allerede
bleven saa stort et Gode, at man ikke vil give Slip paa det igen.
Der leves ikke mindre dybt og rigt udenfor Kirken end i den.
Overfor den kirkelige Tro staar der, og vil for Fremtiden stedse
staa en alvorlig, fri og selvbevidst Livsanskuelse.

De Fritænkere, som mene, at de positive Religioners Tid
er forbi, se ikke, at Kirken netop vil vinde ved Religionsfriheden.
Den Strid, der er rejst mod Kirken i de sidste Aarhundreder,
gælder for største Delen langt mere Kirkens verdslige Magt og
ydre Herredømme end Dogmerne. Striden har ikke staaet saa
meget om Dogmernes Sandhed, som om deres Ret til at bestemme
Individernes borgerlige og moralske Stilling. De Erfaringer, som
ere gjorte i Nordamerika, i Frankrig efter Revolutionen og i
Tyskland efter 1875, godtgøre, at den kirkelige Interesse og Alvor
tiltager med Adskillelsen fra Staten. Hvad der tilsyneladende
svækker Kirken, styrker den i Virkeligheden. De frie Kræfter
vækkes til Virksomhed. Selv den engelske Statskirke, den stiveste
af alle protestantiske Statskirker, bygger mere og mere paa Frivillighedens
Princip. Tocqueville udledte Religionens store Magt i
Nordamerika deraf, at den har trukket sig helt bort fra Politiken.
% --- p. 452 ---
\setcounter{footnote}{0}%
\opage{452}Den har der, i Frihedens forjættede Land, saa stor Indflydelse, at
den offentlige Mening sætter langt snævrere Grænser for fri
Forsken og frie Anskuelser, end den gør her i Europa. Fritænkere,
som tillige ere Demokrater, og som arbejde paa Kirkens Adskillelse
fra Staten, maa derfor have en stor og frejdig Tro paa
Sandhedens Magt. --- Hvad var ogsaa Livet uden en saadan Tro?
Men man kan have den og dog mene, at der er saare langt
til Enden.

Da det er af Betydning for et Spørgsmaal, hvorpaa jeg
senere vil henlede Opmærksomheden, vil jeg her fra et rent
psykologisk Standpunkt fremføre nogle af de Aarsager, som efter
min Opfattelse vedvarende ville medføre Tilslutning til positiv
religiøs Tro. Der er naturligvis saare mange forskellige psykologiske
Forhold, som virke med her\footnote[1]{I min Afhandling „Forholdet mellem Tro og Viden i dets historiske Udvikling“ (Vor Ungdom. 1885) har jeg haft Lejlighed til at omtale nogle Aarsager, som gjorde sig gældende ved den religiøse Reaktion i Begyndelsen af dette Aarhundrede (Se især p. 189---196).}; men der er især to Punkter,
som min Erfaring og mine Undersøgelser have ført mig til at
lægge Vægt paa.

Jo mere den aandelige Udvikling skrider frem, des mere
forskelligartede blive de enkelte Sider eller Retninger af det
aandelige Liv. Teoretisk Forsken og praktisk Interesse, Forstand
og Følelse, Videnskab og Kunst sondre sig mere eller mindre ud
fra hinanden. Hver Side af Livet stiller sine Opgaver paa sin
Maade og løser dem ved sine ejendommelige Metoder. Følgen
bliver let en Delthed af Kraft og Interesse, en Mangel paa Koncentration,
som er uudholdelig for mange Naturer. Der vil opstaa
en naturlig Trang til at samle det spredte paa ét Punkt, eller dog
at have én Region, hvor Sondringen ikke griber ind, én Sfære,
hvor man kan føle det udelte Liv i dets hele Fylde. Der er nu
sikkert flere forskellige Maader, paa hvilke denne Trang vil kunne
tilfredsstilles. Men saa meget er vist, at i de store Folkereligioner
er der saare meget, som imødekommer den. Disse Religioner
have --- hver til sin Tid --- ydet fuld Tilfredsstillelse for alle
Retninger af det aandelige Liv. De skelne ikke imellem det
sande, det skønne og det gode. De henvende sig ikke særlig til
Forstanden eller til Fantasien eller til Følelsen eller til Viljen,
men tale til alle Ævner paa én Gang og forstaas kun, naar de
% --- p. 453 ---
\setcounter{footnote}{0}%
\opage{453}paa én Gang kunne faa dem alle i Tale. Den positive Religion i
sin klassiske Form kender ingen særlig Videnskab, ingen
særlig Kunst, ingen særlig Etik eller Politik. Det betegner altid et
begyndende Brud, naar en af disse Bestræbelser udsondrer sig
som noget for sig og vil gaa sine egne Veje. Ti naar f. Eks.
Videnskaben bliver et eget Omraade ved Siden af Religionen, saa
bliver jo Religionen ogsaa et eget Omraade ved Siden af Videnskaben
og ikke længere Form og Bærer for Livet i dets Helhed.
De store positive Religioner ere ejendommelige Fortætninger eller
Synteser af Tanker og Idealer, Forhaabninger og Bestræbelser,
som af den kritiske Forsken henføres til forskellige Omraader og
under forskellige Synspunkter. Der vil nu være mange, hvis
Natur er saaledes beskaffen, at deres dybest gaaende intellektuelle,
æstetiske og etiske Trang fordrer Tilfredsstillelse under ét, --- for
hvem en Tanke, der ikke tillige tilfredsstiller Fantasiens Trang til
anskuelige Billeder og Følelsens Trang til Haab og Trøst, heller
ikke er nogen \emph{sand} Tanke. De kunne og ville ikke skelne (som vi i
Psykologien skelne) mellem selve Følelsen og de mere eller mindre
symbolske Former, i hvilke den lægger sig for Dagen, eller
mellem en Tankes ideale Værdi og reale Gyldighed. En religiøs
Reaktion kommer ofte efter en følelsesmat og fantasifattig Tid, og
kan da tjæne til Berigelse og Styrkelse for Aandslivet. Det aandelige
Liv har, som alt Liv, en rytmisk Karakter, bevæger sig fremad
gennem skiftende Perioder af Samlen og Spreden, Enhed og Sondring,
Syntese og Analyse. Og der vil stedse være Naturer, som
ere særlig anlagte for den ene af disse Retninger.

Det andet Punkt, som jeg vil fremhæve, staar tilsyneladende
i Strid med det første, men ogsaa kun tilsyneladende. Det er den
Maade, paa hvilken Livets største og stærkeste Modsætninger
komme til Udbrud og Udtryk i de positive Religioner. Den religiøse
Følelse indeholder en Trang til udtømmende at erfare den
største Lidelse og den største Lyst, det største Fald og den største
Oprejsning, som kan tænkes. Kun da mener den at kende Livet,
naar den kender de største Modsætninger, Livet kan medføre; og
kun naar den paa denne Maade kender Livet, stoler den paa at
kunne overvinde Verden. I de store Folkereligioner vækkes de
mest stridende Følelser. Menneskenes Sind føres igennem dem,
ligesom Dante førtes gennem Helvede, Skærsilden og Paradiset.
Den religiøse Kultus virker særlig gennem Kontraster. Ved de
eleusinske Mysterier synes Deltagerne at have fejret Demeters og
% --- p. 454 ---
\setcounter{footnote}{0}%
\opage{454}Persefones Historie: Moderens Sorg ved Datterens Forsvinden,
hendes Sjæleangest og hendes rastløse Søgen, endelig hendes Jubel
ved Genfindelsen. Ogsaa Grækerne havde altsaa deres mater
dolorosa. I den kristne Kirkes Aar gennemleves Kristi Liv paany
af de troende, og dermed en Række af de største Modsætninger
paa Følelseslivets Omraade. --- Der gives Naturer, som kun indenfor
en positiv Religion faa fuld Næring for deres Trang til at
gennemleve Modsætningerne. Der er noget i denne Trang, som
er almenmenneskeligt, og som lægger sig for Dagen hos ethvert
alvorligt Menneske. Den æstetiske Virkning af det tragiske hænger
sammen hermed uden dog ganske at falde sammen med hvad
Talen her er om. Naturer, hos hvem denne Trang gør sig
gældende med særlig Styrke, ville ikke let skelne, som andre gøre,
mellem Følelsen selv og de Symboler, i hvilke den skaffer sig Udtryk
og finder Hvile; Symbolet bliver for dem til Dogme, til det
fuldkomne Udtryk for den højeste Sandhed.

Jeg omtaler dette ogsaa for at pege hen paa, hvor lidet
den blotte logiske og historiske Kritik paa det religiøse Omraade
trænger ind til Sagens Kærne. Bag de logiske Modsigelser og
historiske Usandsynligheder, som Kritiken formaar at paavise,
ligge de mægtige Følelsesbevægelser og Følelsesmodsætninger.
Det er disse, der gør Tilværelsen til saa stor en Gaade; hvor kan
den rumme de Modsætninger, som fremkalde disse stridende
Følelser? hvorledes kunne de udspringe fra én og samme Kilde?
hvorledes kunne de forsones og overvindes? --- Det er klart, at man
meget vel kan anerkende og føle Modsætningernes Virkelighed
uden at godkende Forsøgene paa at løse hin Gaade. Man kan erkende,
at Tilværelsen viser os baade den største Lidelse og den
største Lykke, baade den største Fornedrelse og den mest ophøjede
Storhed, og dog ikke forstaa, hvorledes de kunne skyldes
samme Aarsag, især naar denne skal være et absolut godt Væsen.
Dogmet staar her --- efter min Overbevisning --- lige saa afmægtigt
som Videnskaben. Men den fuldt berettigede Kritik af Løsningerne
bør ikke bevirke Miskendelse af de stærke Drifter, der føre til at
stille Spørgsmaalene. Ti derved gør man Livet snævert og fattigt.
Vi staa endnu --- trods al vor Religionsfilosofi og vor sammenlignende
Religionsvidenskab --- meget langt tilbage i Henseende
til Forstaaelse af det psykologiske Grundlag for de positive
Religioner.

% --- p. 455 ---
\setcounter{footnote}{0}%
\opage{455}Jeg er kommen noget bort fra mit Hovedæmne og vender
nu tilbage til det.

\sectionnum{III.}

Saa meget er vist, at Kampen  vil blive lang. Den vil efter
manges Mening blive den evige Strid. Skulde det være saa, maa
vi gøre os fortrolige dermed. For en alvorlig Livsanskuelse er
hele Livet en stor Kamp, og denne vil da kun være en af dens
mange Former. Men det er ikke ligegyldigt, paa hvilket Grundlag
og paa hvilken Maade Striden føres. Skal den fra kirkelig
Side føres med det Skriftsted som Motto, med hvilket en amerikansk
Præst rykkede ud mod Darwin og Spencer: »Skulde jeg ikke hade
dem, der hade dig, Herre!«? Og skal den fra den modsatte Side
føres med Haan og Spot? --- For de værste Udskejelser i denne
Retning er der jo sørget ved Religionsfrihedens officielle Anerkendelse.
Det ydre, borgerlige Liv kan gaa sin Gang, uanfægtet
af de enkelte Individers forskellige Tro. De fanatiske troende og
ikke-troende maa i hvert Tilfælde taale hverandres Nærværelse i
Samfundet. Men det er, som allerede tidligere bemærket, et vel
udvortes og negativt Forhold.

Det vil heller ikke være nok at henvise til den stedse mere
udbredte og stedse bedre begrundede Overbevisning, at de religiøse
Problemer angaa noget, som i Følge sin Natur ligger udenfor
den strænge Videnskabs Rækkevidde. Ti Spørgsmaalet er netop:
hvad saa? Skulle vi lade hverandre hver sejle sin egen Sø, fordi
man dog intet sikkert kan vide om disse Ting? Vil det ikke føre
til Ligegyldighed eller til Afbrydelse af enhver dybere gaaende
aandelig Forbindelse mellem dem, der ikke høre til samme Sekt?
--- Og i hvert Tilfælde vilde man da ikke kunne fordre Tolerance
af dem, der erklære sig at være aldeles sikre i deres Sag, ---
sikrere end nogen som helst Videnskab kunde gøre dem. Katolikerne
paastaa saaledes ofte\footnote[1]{Se f. Eks. en Udtalelse af en katolsk Biskop, som anføres hos F. C. Baur: Kirchengeschichte des 19ten Jahrhunderts. P. 264.}, at Protestanterne \emph{maa være} tolerante,
fordi de ikke tør hævde, at deres Kirke er den ene sande,
men at Katolikerne ikke \emph{kunne være} det, fordi det er en Del af
deres Tro, at deres Kirke er den ene sande. Protestanterne lade
maaske Paastanden gaa videre til Fritænkerne --- men de Fritænkere,
der mene, at alle troende enten maa være dumme eller
% --- p. 456 ---
\setcounter{footnote}{0}%
\opage{456}hykleriske, siden de kunne tro, hvad de tro, ville sende den tilbage
med Protest.

Skal Religionsfriheden have mere end ydre og negativ Betydning,
maa vi undersøge, om der ikke skulde være fælles Omraader,
som kunne unddrages Striden, skønt de ere af væsentlig
Indflydelse for det menneskelige Liv, og om Muligheden af gensidig
Forstaaelse og Sympati er aldeles afskaaren for selve de Spørgsmaals
Vedkommende, om hvilke Striden staar. I første Henseende
vilde det især være af stor Betydning, om der gives en human,
paa naturligt menneskeligt Grundlag byggende Etik. I sidste Henseende
vil det komme an paa, om Mennesker, der staa paa forskellige
religiøse Standpunkter, alligevel kunne have et vist aandeligt
Slægtskab med hverandre.

Dersom det eneste Synspunkt, de forskellige religiøse Standpunkter
kende overfor hverandre, er: Sandhed eller Løgn, saa
nytter det ikke at undersøge disse Spørgsmaal. Men i saa Fald
er det heller ikke godt at se, hvorledes man kan gaa ind paa den
ydre, borgerlige Religionsfrihed. Hvor kan man finde sig i, at
Løgnere udøve Borgerrettigheder og deltage i de offentlige Anliggenders
Afgørelse? Og det er heller ikke godt at se, hvorledes
man kan anerkende Videnskabens Selvstændighed. Gives der en
Videnskab, som naa'r sit Maal uden Hensyn til og uden Ledelse
af religiøse Anskuelser, saa kan der dog ad menneskelig Vej naa's
i hvert Tilfælde til en Del af Sandheden. Ti hvorledes man saa
dømmer Viden og Tro imellem, --- tilsidst kan der dog kun være
én Sandhed. --- Men skulde den menneskelige Naturs Ævne til at
udvikle sig ved sine egne Kræfter være indskrænket til det borgerlige
Liv og til Videnskaben? Gives der ikke ogsaa en rent
menneskelig Etik, som, i hvert Tilfælde til en vis Grad, er uafhængig
af religiøse Antagelser? --- Jeg forstaar ikke, hvorledes
nogen, der i ethvert Menneske ser et ejendommeligt personligt
Væsen, kan nægte dette. Et personligt Væsen kan kun i sig selv 
have den sidste Maalestok for, hvad det skal anerkende som sandt,
ædelt og godt. Hvad enten det er noget overnaturligt eller noget
naturligt, Personligheden staar overfor: i sit eget Indre maa den
søge Principerne for Forstaaelse og Vurdering af alt hvad der
fremstiller sig for den. Hvilke Idealer Mennesket saa anerkender,
og hvor ophøjet endog den etiske Lov staar \emph{for ham}, de skulle
jo dog være Ideal og Lov for ham, og blive det kun ved at
finde Tilknytningspunkter i hans indre aandelige Natur. Det er
% --- p. 457 ---
\setcounter{footnote}{0}%
\opage{457}den store Betydning af Sokrates's \emph{Bud: Kend dig selv}! ved
hvilket han blev Etikens Grundlægger. Den filosofiske Etik har i
sine forskellige Former bygget paa dette Princip, som kan kaldes
\emph{Personlighedsprincipet}, Principet baade for den frie Forsken
og den frie Samvittighed. Idet den filosofiske Etik arbejder med
dette Princip, drager den derved tillige Konsekvenserne af den
almindelige Menneskekærligheds Princip. Den kan ikke gaa ind
paa den Betragtningsmaade, at Forskellen mellem Antagelse eller
Ikke-Antagelse af visse religiøse Lærdomme skal drage en absolut
Grænse mellem Menneskene, evig skille dem i to Dele. Som Menneskekærligheden
har gennembrudt de Skranker, der skilte Nation
fra Nation og Stand fra Stand, vil og maa den ogsaa gennembryde
de Skranker, som konfessionelle Forskelligheder have opstillet
for den fulde Anerkendelse af det personlige Liv i alle
menneskelige Væsener. Den religiøse Strid vil ikke dermed ophøre,
og den vil ikke miste sin Skarphed. Der skal intet slaas af
i Spørgsmaalet om, hvad Sandhed er. Men Striden kan være
skarp uden at være bitter og hadsk.

Dersom de troende kun kunne betragte de ikke-troendes
Dyder som glimrende Laster, --- dersom Ikke-Antagelsen af Kirkens
Lære kun er at udlede af ond Vilje og Antipati mod Helligheden,
--- dersom Striden mellem Tro og Ikke-Tro fuldstændig er ét med
Striden mellem godt og ondt, saa er alt Haab ude. Ofte lyder
det, som om det var saa. Men meget af det, der lyder, er ogsaa
kun store Ord, ikke Tanker, i hvert Tilfælde ikke Tanker, hvis Konsekvenser
man drager. Ti den fulde Konsekvens af saadanne
Tanker vilde føre langt tilbage --- tilbage til Middelalderen, til
Tiden, som ligger forud for Religionsfrihedens Anerkendelse og
den moderne Kulturs Udvikling!

For den troende vil naturligvis den rent menneskelige Etik
ikke kunne være udtømmende. Men det maa dog være af Interesse
ogsaa for ham at se, hvor langt man ad denne Vej kan
naa, og han vil meget vel selv kunne deltage i Undersøgelsen.
Han vil forbeholde sig at sætte sit eget Spir paa Bygningen.
Men der vil i hvert Tilfælde være et fredlyst Omraade af stor Betydning,
hvor den religiøse Kamp ikke behøver at trænge ind, lige saa
lidt som den nu trænger ind paa Fysikens og Kemiens Enemærker\footnote[1]{Angaaende den humane Etiks Væsen og Betydning tillader jeg mig at henvise dels til mit Skrift „Om Grundlaget for den humane Etik“ (Køben%
\opage{458}havn 1876), dels til min Afhandling „Om den filosofiske Etiks Principer“ i „Oversigt over det kgl. danske Videnskabernes Selskabs Forhandlinger i Aaret 1886“.}.
% the note runs over from p. 3457 to p. 3458 (joined by draft.py)

% --- p. 458 ---
\setcounter{footnote}{0}%
\opage{458}Først naar det anerkendes, at Etiken har sit Grundlag i
selve den menneskelige Natur, vil den egentlige, positive Religionsfrihed
blive mulig, idet den største Hindring er bortryddet for
gensidig Forstaaelse og Sympati mellem Mennesker, som staa paa
forskelligt religiøst Standpunkt.

Der er dog endnu en vigtig Betingelse, som maa være tilstede,
og som --- til Dels formedelst den dogmatiske Maade, paa
hvilken den religiøse Kamp, baade som Angreb og som Forsvar,
er bleven ført --- endnu kun i ringe Grad er udviklet iblandt os.
Det er Sansen for det personlige Liv, særlig Følelseslivet, under
alle dets forskellige Former. Den hænger ogsaa meget nøje
sammen med Troen paa, at der i ethvert Menneske rører sig, eller
i ethvert Tilfælde er Anlæg til et ejendommeligt Liv. Den forudsætter
en Glæde ved at høre hver Fugl synge med sit Næb. Men
den kræver en særegen Opdragelse. Vi ere endnu altfor dogmatiske
i vor Psykologi. Vi gaa altfor trygt ud fra, at andre
opfatte, tænke og føle ganske paa samme Maade som vi. Der gør
sig individuelle Forskelligheder gældende selv ved de simpleste
aandelige Virksomheder, Forskelligheder, som man kun lige har
begyndt paa at drage frem for Lyset. Hvor meget mere maa da
det højere aandelige Liv forme sig forskellig, og hvor forsigtige
maa vi være, for ikke at drage forhastede Slutninger fra den
Maade, hvorpaa Følelseslivet giver sig Udtryk i Ord eller Symboler!
Disse kunne lyde ens og dog dække over indre Bevægelser
af meget forskellig Art. Kun den, der har Ævne til at skelne
mellem Skal og Kærne, opdager disse Forskelligheder, denne
Rigdom. Den, der standser ved den positive eller negative Trosbekendelse,
kommer ikke længere end til Overfladen. Blomsten,
som vokser i det høje Græs, opdages kun af den, der bøjer de
lange Straa til Side. Men den, der gør det, vil ikke blot finde
store Forskelligheder og rige Nuancer, hvor han ikke anede det.
Han vil ogsaa opdage, at det aandelige Slægtskab er et andet
end det dogmatiske. Ligesom der bag samme Trosbekendelse
kan ligge et højst forskelligt indre Følelsesliv, saaledes kan der
bag forskellig Trosbekendelse ligge et nøje beslægtet Følelsesliv.
Det kan være vanskeligt at trænge ind dertil, og vi staa langt
% --- p. 459 ---
\setcounter{footnote}{0}%
\opage{459}tilbage i denne Kunst. Men før man lærer det, har man endnu
ikke lært det væsentlige i Mennesket at kende. Som allerede
Schiller har sagt; hvad vi tænke, tilhører alle; kun det, vi føle,
er vort eget. Og Følelseslivet har stedse en vis Uafhængighed af
Tanker og Forestillinger. Lad os tage et enkelt Eksempel. For
den alvorlige troende som for den alvorlige Fritænker staar
Livet som en Kamp, i hvilken det gælder at hævde Idealet; de
føle det tragiske i denne Kamp og de store Modsætninger, gennem
hvilke den fører; deres Livsfølelse, deres Følelse ved Tilværelsen
er tilsidst bestemt ved disse Erfaringer. Skulde der ikke være
større aandeligt Slægtskab mellem to saadanne Mennesker end
mellem saadanne, der høre til samme Konfession, men hvis Følelsesliv
bevæger sig i helt forskellige Baner? --- Psykologisk set er Forskellen
mellem den troende og Fritænkeren ofte kun den, at hin
har sine dogmatiske Former, i hvilke han finder Udtryk for sine
inderste Livserfaringer, medens Fritænkeren maa sige med Faust:
\begin{verse}\small
Ich habe keinen Namen dafür!\\
Gefühl ich alles,\\ % PRINTED AS IS (ich: Goethe: „Gefühl ist alles“; zoom z1, reader K)
Name ist Schall und Rauch.
\end{verse}

Under den staaende, maaske evige Strid vilde det være et
stort Gode, om en Forstaaelse af den Art, jeg her har antydet,
kunde udvikle sig. Den vilde ikke blot være en Berigelse og Udvidelse
af vort Liv, den vilde være den eneste mulige Form for
Menneskeslægtens aandelige Enhed, efter at den store Tanke om
at forene alle under én Bekendelse har maattet bukke under. Religionsfriheden
vilde da ikke blot være noget ydre og negativt,
men uden at slaa noget af paa hvad den store Strid væsentlig
drejer sig om, vilde de samlende og forenende Kræfter kunne virke.
\signature{\textsc{Harald Høffding}.}

\end{document}
